\section{Step 39 Explain a block boundary in the measured work}\label{pcm-s39}

Compare complete calls rather than only the addition instructions. At sixty-three samples, the MSA path executes seven full groups and seven scalar samples. At sixty-four it executes eight groups. The additional sample can therefore reduce executed work. Include dispatch and tail costs when explaining this result, and keep simulated instruction counts distinct from host time.

\subsection{Expected observation}
Quiet sixty-three-sample calls use 359 MSA-path instructions; sixty-four use 169. The corresponding scalar counts are 1896 and 1926.

\Needspace{11\baselineskip}
\subsection{Check your understanding}
Why can the sixty-four-sample MSA call execute fewer instructions?

\begin{itemize}
\item \textbf{A} A full group replaces the seven-sample scalar tail.
\item \textbf{B} The last sample is skipped.
\item \textbf{C} The elapsed time is defined by sample count alone.
\end{itemize}
